\Askhsh{34813}{2}
\begin{ylh}{1}
	\item 2.2 Παρουσίαση Στατιστικών Δεδομένων
\end{ylh}
Ο αριθμός των ετήσιων επισκέψεων (άξονας των $x$) ενός δείγματος 160 μαθητών ενός σχολείου στα διάφορα μουσεία της χώρας δίνεται από το παρακάτω διάγραμμα σχετικών συχνοτήτων.

\begin{center}
\begin{tikzpicture}
\begin{axis}[
    axis x line=bottom,
    axis y line=left,
    axis line style={->},
    xlabel = {$x$},
    ylabel = {$f_i$},
    xmin = -0.5, xmax = 4.5,
    ymin = 0, ymax = 0.35,
    xtick = {0,1,2,3,4},
    ytick = {0.1, 0.15, 0.2, 0.25, 0.3},
    yticklabels = {0,1, 0,15, 0,2, 0,25, 0,3},
    yticklabel style={font=\small, anchor=east},
    xticklabel style={font=\small, yshift=-3pt},
    grid style = {dotted},
    width = 9cm, height=7cm,
    every axis x label/.style={at={(current axis.right of origin)}, anchor=west},
    every axis y label/.style={at={(current axis.above origin)}, anchor=south, rotate=0},
    clip=false
]
\addplot[ycomb, black, mark=*, mark size=1.2pt, thick] coordinates {
    (0, 0.2)
    (1, 0.3)
    (2, 0.25)
    (3, 0.15)
    (4, 0.1)
};
\end{axis}
\end{tikzpicture}
\end{center}

Για το παραπάνω δείγμα βρείτε:
\begin{alist}
	\item το ποσοστό των μαθητών που έκανε ακριβώς δύο επισκέψεις ετησίως. \monades{8}
	\item πόσοι μαθητές πραγματοποίησαν ακριβώς μία επίσκεψη ετησίως. \monades{8}
	\item το ποσοστό των μαθητών που έκανε δύο τουλάχιστον επισκέψεις ετησίως. \monades{9}
\end{alist}

